\documentclass[11pt]{article}

\usepackage[french, english]{babel}

\usepackage[T1]{fontenc}

\usepackage{amsfonts}
\usepackage{amsmath}
\usepackage{amssymb}

\usepackage{graphicx}
\usepackage{multirow}
\usepackage{multicol}
\usepackage{xcolor}
\usepackage[shortlabels]{enumitem}

\usepackage[margin=2cm]{geometry}

\usepackage{mathtools}
\usepackage{siunitx}
\usepackage{braket}

\usepackage{listings}

\begin{document}

\begin{multicols}{2}
    \section{\underline{Newton}}
    \subsection{Mvt circulaire}
    Position : pointe l'emplacement de la particule dans le système.\\
    Vitesse : Tangantielle à la trajectoire.\\
    Accélération : Pointe vers le centre du cercle de mouvement.\\

    \subsection{Mvt parabolique}
    \begin{equation}
        \vec{r}=\vec{r}_0+\vec{v}_0t + \frac{\vec{a}t^2}{2}
    \end{equation}

    
\end{multicols}
\end{document}
